%! Periodic Phenomena — Unit 6, Lesson 6.1 — Homework (Answer Key)
\documentclass[10pt]{article}
\usepackage{precalculus-article}
\usepackage{precalculus-key}   % pulls in -boxes; do NOT also load -boxes
\newcommand{\TallMath}[1]{$\displaystyle #1\rule[-1.4em]{0pt}{3.2em}$}

\begin{document}

\pageheader{Unit 6, Lesson 6.1}{Homework --- Answer Key}
\namedateperiod

\vspace{0.10in}

\begin{notesbox}{Problems 1--6 --- Key}

\textbf{1.}\quad

\vspace{4pt}
\textbf{(a)}
\begin{center}
\renewcommand{\arraystretch}{1.35}
\begin{tabular}{|c|c|c|c|c|c|c|c|c|}
\hline
$x$ & 0 & 2 & 4 & 6 & 8 & 10 & 12 & 14 \\
\hline
$f(x)$ & 1 & 4 & 7 & 4 & 1 & 4 & 7 & 4 \\
\hline
\end{tabular}
\end{center}
Periodic? \ans{Yes} \quad Period $k =$ \ans{8} \quad
Reason: \ans{The sequence 1, 4, 7, 4 repeats every 8 units of input (e.g., $f(0)=f(8)=1$, $f(2)=f(10)=4$, etc.). The smallest such input length is 8.}

\vspace{8pt}
\textbf{(b)}
\begin{center}
\renewcommand{\arraystretch}{1.35}
\begin{tabular}{|c|c|c|c|c|c|c|c|c|}
\hline
$x$ & 0 & 1 & 2 & 3 & 4 & 5 & 6 & 7 \\
\hline
$g(x)$ & 2 & 5 & 9 & 14 & 20 & 27 & 35 & 44 \\
\hline
\end{tabular}
\end{center}
Periodic? \ans{No.} \quad Reason: \ans{The outputs increase without bound and never repeat. The differences between consecutive outputs are 3, 4, 5, 6, 7, \ldots, growing each time. A periodic function must return to the same output values repeatedly.}

\vspace{0.12in}
\textbf{2.}\quad $h$ has period $k = 5$; $h(2) = 13$, $h(3) = -1$.

(a)\quad $h(7) = $ \ans{$h(2+5) = h(2) = 13$. Because $h$ has period 5, $h(7) = h(2) = 13$.}

(b)\quad $h(12) = $ \ans{$h(2 + 2\cdot 5) = h(2) = 13$. Every 5 units the output repeats, so $h(12) = 13$.}

(c)\quad \ans{$h(x) = -1$ whenever $x = 3 + 5n$ for any integer $n$: that is, $x = \ldots, -7, -2, 3, 8, 13, 18, \ldots$}

\vspace{0.12in}
\textbf{3.}\quad Tidal depth.

(a)\quad Period: \ans{12 hours (the pattern restarts at $t = 12$, identical to $t = 0$).}

(b)\quad Maximum depth: \ans{8 m, occurring at $t = 4$ hours.}

(c)\quad \ans{The depth equals 6 m at approximately $t \approx 3$ hr and $t \approx 5$ hr (between the 5.5 and 8 readings). The depth exceeds 6 m on approximately $3 \leq t \leq 5$ hours. (Accept any reasonable interpolation.)}

(d)\quad \ans{Within the first 12-hour period, depth $= 8$ m at $t = 4$ hr. The pattern repeats every 12 hours, so depth $= 8$ m also at $t = 16$ hr. (Two times in 24 hours.)}

\vspace{0.12in}
\textbf{4. (AP-Style MC)}\quad

Answer: \ans{(C)} \quad \ans{By definition, a periodic function with period 6 satisfies $p(x+6)=p(x)$ for all $x$ in the domain. Choice (A) would require $p(0)=0$ as well, which is not given. (B) is incorrect --- the output repeats, it does not increase by 6. (D) and (E) are not required by the definition of periodic.}

\vspace{0.12in}
\textbf{5. (AP-Style MC)}\quad

Answer: \ans{(C)} \quad \ans{The sequence of outputs is 6, 9, 6, 3, 6, 9, 6, 3, 6, 9, \ldots The pattern 6, 9, 6, 3 repeats every 4 units of input, so the period is 4.}

\end{notesbox}

\vspace{0.08in}

\begin{extensionbox}
\textbf{Extension (Key):}

(a)\quad \ans{The average rate of change of $s$ on $[12, 16]$ is also 3. Since the period is 10, the interval $[12, 16]$ corresponds to $[2, 6]$ shifted by one period ($12 = 2 + 10$, $16 = 6 + 10$). Because $s(x+10)=s(x)$, the output values are identical, so all average rates of change over corresponding intervals are equal.}

(b)\quad \ans{Not necessarily. The average rate of change on $[0,10]$ is $\dfrac{s(10)-s(0)}{10}$. Since $s$ has period 10, $s(10) = s(0)$, so the average rate of change over a full period is $\dfrac{0}{10} = 0$. Even though $s$ is increasing on part of the period, it must decrease just as much to return to the starting value.}
\end{extensionbox}

\begin{teachernote}
\textbf{Problem 2 guidance:} Students often want to compute $h(7)$ by adding
something rather than recognizing that the period means the output \emph{repeats
identically}. Prompt: ``What does $h(2+5)$ equal by definition of period?''

\textbf{Problem 5 MC:} Students may incorrectly choose period 2 (noticing
6 appears at $x=0,2,4,6,\ldots$). Point out that 6 appearing every 2 units
does not mean the pattern repeats every 2 units --- check all outputs. The
full repeating block is 6, 9, 6, 3 which has length 4.

\textbf{Extension (b):} This is a key conceptual point --- the average rate
over a full period is always 0 for any periodic function, because the output
returns to its starting value. Connects to the Fundamental Theorem of Calculus
in future courses.
\end{teachernote}

\begin{tcolorbox}[colback=navylight, colframe=navy, arc=2mm,
    left=3mm, right=3mm, top=2mm, bottom=2mm]
\small\textbf{Preview --- Lesson 3.2: Sine, Cosine, and Tangent}\\
We've been identifying and describing periodic relationships from tables.
In Lesson 3.2, we'll meet two functions --- sine and cosine --- that are
the \emph{mathematical model} for the smoothly repeating patterns we saw
today. Come ready to connect period, increase/decrease, and concavity
to the vocabulary of sinusoidal functions.
\end{tcolorbox}

\end{document}
